% A.102: a full-congruence-plus-abelian quotient with all bad-place comparisons.
\begin{statement}{A.102}{Proved}
{All bad-place comparisons can hold while a good place obstructs the quotient section}
Put $\lambda=254$, let $s_\lambda$ be its rational point
section, and write
\[
 C_H=\ker(\Pi\longrightarrow H),\qquad
 \mathcal Q=E_U/(C_H\cap\Gamma_2\Pi).
\tag{A.102.1}
\]
There is an explicit continuous section
$\sigma:G_\Q\to\mathcal Q$ with the following properties.
Its projection to $E_H$ is literally the projection of
$s_\lambda$.  Its completed Kummer coordinates have finite
support contained in $\{2,11,23,127\}$ and are different
from $(254,-253)$.

At the four bad primes of $E_{254}$ its local section
classes in this same quotient are as follows:
\[
\begin{array}{c|c}
 \text{physical prime}&\text{actual local source of the quotient section}\\ \hline
 2,11&\text{the point }254\\
 127&\text{a full local cusp-0 section}\\
 23&\text{a full local cusp-1 section}.
\end{array}
\tag{A.102.2}
\]
These are simultaneous comparisons in $\mathcal Q$,
with geometric conjugators, rather than separate
comparisons in its two factors.  The Tate comparisons
preserve the integral Weil pairings.

At the good prime $47$, the section is the image of
neither an ordinary point nor any of the three cusp
packets.  The obstruction is already visible in the
finite geometric abelian quotient at exponent $46$.
Consequently $\sigma$ is not the image of a Selmer section
of $E_U$.  No lift to a full section is constructed.
This gives an exact example in which the global full-$H$
comparison, finite Kummer support, and every bad-place
point or cusp comparison all hold, while a good-place
comparison fails.

Moreover, every noncuspidal Selmer section in A.14 with
arbitrary labeled integral realization $E_{254}$ is
$s_{254}$.  The proof uses the forced coefficient-$23$
data from the surviving bad-place configuration in
the residue field at $47$.  This conclusion requires
no full-$H$ hypothesis on that section.
\end{statement}
\begin{proof}
\emph{The four explicit local exponents.}
We use the branch of $\log_p$ with $\log_p(p)=0$ and
with roots of unity in its kernel.
For $p=127,23$, let $q_p\in\Q_p$ be the Tate nome
determined by $j(E_{254})$; it has valuation two.
It is the unique small solution of
$j(q_p)=j(E_{254})$.
The exact inverse-nome expansion of A.74 gives
\[
 \begin{split}
       q_{127}/127^2&=2^{-6}F(254),\\
       q_{23}/23^2&=11^2\,2^{-8}F(-253),\\
       F(w)&=1+w+O(w^2),\qquad
                     F(w)\in1+w\Z_p[[w]].
 \end{split}
\tag{A.102.3}
\]
Here the last notation is evaluated at the respective
$w\in p\Z_p$; the power series has coefficients integral
away from two.  The equality of the two Legendre
$j$-expressions at $t$ and $1-t$ justifies using
$w=-253$ at the prime $23$.

Define four actual $p$-adic numbers by
\[
 \begin{array}{ll}
 a_{127}=\displaystyle
       \frac{-6\log_{127}(2)}{\log_{127}(q_{127}/127^2)},
 &b_{127}=\displaystyle
       \frac{-\log_{127}(11)}{\log_{127}(23)},\\[8pt]
 a_{23}=\displaystyle
       \frac{-\log_{23}(2)}{\log_{23}(127)},
 &b_{23}=\displaystyle
       \frac{2\log_{23}(11)-8\log_{23}(2)}
                         {\log_{23}(q_{23}/23^2)}.
 \end{array}
\tag{A.102.4}
\]
Every numerator and denominator here has valuation one.
For a $p$-adic unit $r$, its exact first logarithm residue is
\[
       \frac{\log_p(r)}p\equiv
          \frac{r^{p-1}-1}{p(p-1)}\pmod p.
\tag{A.102.5}
\]
Using this finite integer formula and (A.102.3) gives
\[
\begin{array}{c|rrrr}
 p&\log_p(2)/p&\log_p(11)/p&
            \log_p(r)/p&\log_p(q_p/p^2)/p\\ \hline
 127&109&35&90\ (r=23)&110\\
 23&6&16&5\ (r=127)&19.
\end{array}
\tag{A.102.6}
\]
It follows that
\[
        (a_{127},b_{127})\equiv(31,49)\pmod{127},
          \qquad (a_{23},b_{23})\equiv(8,4)\pmod{23}.
\tag{A.102.7}
\]
The displayed residues have respectively the square roots
$44,7$ modulo $127$ and $10,2$ modulo $23$.
Their derivatives under squaring are units, so Hensel's
lemma proves that all four numbers in (A.102.4) are
squares of $p$-adic units.  In particular they are
well-defined units, and none equals one.

Define $\alpha,\beta\in\widehat\Z^\times$ by their
actual primary coordinates:
\[
 \begin{split}
  &\alpha_{127}=a_{127},\quad\alpha_{23}=a_{23},
       \quad \alpha_\ell=1\quad(\ell\ne23,127),\\
  &\beta_{127}=b_{127},\quad\beta_{23}=b_{23},
       \quad \beta_\ell=1\quad(\ell\ne23,127).
 \end{split}
\tag{A.102.8}
\]
Both are squares of units in $\widehat\Z$.
Set
\[
                 A=2\,127^\alpha,
                  \qquad B=-11\,23^\beta
                           \quad\text{in }\widehat{\Q^\times}.
\tag{A.102.9}
\]
Exponentiation here is the continuous map from
$\widehat\Z$ to the power completion obtained by
exponentiating the indicated rational prime in each
finite quotient.  Thus these are actual completed
classes with the asserted finite support.
Their valuation at $127$, respectively $23$, differs
from that of $254$, respectively $-253$, by (A.102.7).

\emph{The global quotient section and its exact product structure.}
Use the point $254$ to split all arithmetic extensions.
By A.56 the actual geometric kernel of $\mathcal Q$ is
\[
          Q=H\mathbin{\times}_{H^{\rm ab}}T,
              \qquad T=\widehat\Z(1)^2.
\tag{A.102.10}
\]
For $\ell=23,127$ choose continuous global Kummer
cocycles for the rational primes $127$ and $23$, and
multiply them by $\alpha_\ell-1$ and $\beta_\ell-1$
in the first and second coordinates of $T_\ell$.
Set every other primary coordinate of this relative
cocycle $z:G_\Q\to T$ equal to zero.
It represents exactly
$(A/254,B/(-253))$.
Since the map $T\to H^{\rm ab}$ has only two- and
three-primary factors, its value on $z$ is literally
zero, not just a coboundary.
Pair the point section in $E_H$ with the abelian
section $z\cdot s_{254}^{\rm ab}$.
The two maps to $E_{H^{\rm ab}}$ agree, so the
topological isomorphism (A.56.4) gives the claimed
continuous section $\sigma$ of $\mathcal Q$.

The product decomposition needed to verify simultaneous
local conjugacy is the following exact one:
\[
 \begin{split}
  Q&=Q_{2,3}\times
          \prod_{\ell\ge5}\bigl(H_\ell\times T_\ell\bigr),\\
  Q_{2,3}&=(H_2\times H_3)
       \mathbin{\times}_{H_2^{\rm ab}\times H_3^{\rm ab}}
                          (T_2\times T_3).
 \end{split}
\tag{A.102.11}
\]
This follows by applying (A.56.1)--(A.56.2) separately
to the primary coordinates.  The point-induced
$G_\Q$-action preserves every factor: it acts on
$H_\ell$ by the point's Tate representation and on
$T_\ell$ cyclotomically.  Thus the corresponding
arithmetic extension is the product of these split
extensions over their common quotient $G_\Q$.
The projection of $\sigma$ to the $Q_{2,3}$ extension
is literally the point section.  On each factor
$H_\ell\times T_\ell$ for $\ell\ge5$, the two
geometric conjugators can be chosen independently.

\emph{The actual local cusp parameters and paired Tate comparisons.}
At the physical prime $127$, the class $B_{127}$ is
one.  Its valuation is zero; its principal-unit
logarithm is
$\log_{127}(11)+b_{127}\log_{127}(23)=0$.
Its root-of-unity coordinate is also one.
Indeed the unmodified value $-253$ is one modulo $127$,
and the only changed coefficient other than $127$ is
$23$, which does not divide $126$.
The decomposition of $\Q_{127}^\times$ into valuation,
roots of unity, and principal units proves the assertion
in the full completion, including every coefficient.
Consequently the pair $(A_{127},B_{127})$ is that of
a full local cusp-0 section with completed tangent
parameter $C=A_{127}$.

Similarly $A_{23}=1$:
$\log_{23}(2)+a_{23}\log_{23}(127)=0$,
the unmodified value $254$ is one modulo $23$, and
the only other changed coefficient $127$ does not
divide $22$.
The pair at $23$ is therefore the pair of a full
local cusp-1 section with tangent parameter $D=B_{23}$.
The cusp decomposition group with its tangent-one
splitting is $\widehat\Z(1)\rtimes G_{\Q_p}$;
the Kummer cocycles of $C,D$ define these full local
sections explicitly.

The exact formal Tate classes from A.63 are
$C^2/256$ at cusp zero and $D^2/256$ at cusp one.
They satisfy full, all-coefficient equalities
\[
       \kappa(C^2/256)=\alpha\,\kappa(q_{127}),
          \qquad
       \kappa(D^2/256)=\beta\,\kappa(q_{23}).
\tag{A.102.12}
\]
For the first equality, both valuation coordinates
are $2\alpha$.  At coefficient $127$, its unit
coordinates agree by the definition of $a_{127}$.
At every other coefficient, $F(254)\in1+127\Z_{127}$
has trivial component.  The only remaining changed
coefficient $23$ sees no unit torsion, since
$23\nmid126$; all the coefficients dividing $126$
retain $\alpha_\ell=1$.  This verifies every component
of the first equality.
For the second, the valuations are $2\beta$, the
coefficient-$23$ unit equality is the definition of
$b_{23}$, and $F(-253)$ is principal at $23$.
The changed coefficient $127$ sees no unit torsion
there because $127\nmid22$, and all other unit
torsion components have $\beta_\ell=1$.

At cusp zero the common quadratic character is
$\chi_{-1}$; at cusp one it is trivial.  These are
also the characters of the respective nodal
Legendre reductions: the double root at zero has
slopes $\pm\sqrt{-1}$, while the double root at
one has slopes $\pm1$.
Cancel the same character on both sides.
In Weil-paired Tate bases, an equality of extension
classes $\kappa(q')=u_\ell^2\kappa(q)$ gives an
intertwiner with diagonal entries
$u_\ell,u_\ell^{-1}$; an upper unipotent change
accounts for any difference of cocycle representatives.
It has determinant one.  Applying this to square
roots of $\alpha_\ell$ or $\beta_\ell$ in
(A.102.12) supplies the paired integral intertwiners
at every coefficient, including $\ell=p$.
For $\ell\ge5$ these matrices lie in
$H_\ell=\operatorname{SL}_2(\Z_\ell)$.
The arithmetic map to
$\operatorname{GL}_2(\Z_\ell)\times G_{\Q_p}$
is injective: its kernel is geometric and the
geometric image $H_\ell$ is faithful by definition.
Hence these are actual $H_\ell$-geometric conjugators
between the corresponding quotient sections.

\emph{Simultaneous conjugacy in the common two- and three-primary factor.}
We give the precise extension of the disc comparison
of A.49 that is needed here.
For a finite extension $K/\Q_p$, put
$\mathcal T=\operatorname{Spec}\mathcal O_K[[w]][1/(\pi w)]$.
For a finite geometric Galois cover whose group has
order divisible only by two and three, with $p>3$,
the normalization argument of A.49 applies with a
$2^a3^b$ common multiple of the boundary ramification
indices.  Its ramification is tame at both divisors:
its degree is prime to $p$ at the vertical divisor,
and the horizontal residue characteristic is zero.
Adjoin the indicated roots of $\pi$ and of $w$.
The normalized base change is finite etale by the
same Abhyankar lemma and purity, and a constant
unramified extension splits it over the complete
local ring.  Thus the geometric pro-$\{2,3\}$ group
of $\mathcal T$ is exactly
$\Z_2(1)\times\Z_3(1)$, with its stated Kummer
coordinate.  The map from $K((w))$ is the identity
on this kernel and on $G_K$, so is an isomorphism
of the corresponding arithmetic extensions.
Evaluation at a nonzero $w\in\mathfrak m_K$
therefore identifies the point with the cusp
section having its two- and three-primary tangent
classes, with one actual geometric conjugator.

The group $Q_{2,3}$ in (A.102.11) is pro-$\{2,3\}$:
$H_2$ is pro-two and $H_3$ has reduction of order
$24$ with pro-three kernel.
The quotient map from $\Pi$ to this factor is
surjective by A.56, so the disc comparison just
proved applies to its entire arithmetic extension.
At $127$ take $w=t$; at $23$ take $w=1-t$.
The point has respective values $254$ and $-253$.
The two- and three-primary components of the
parameters $C,D$ are those same classes, since
$\alpha_2=\alpha_3=\beta_2=\beta_3=1$.
Thus the actual point section and the selected
full cusp section have conjugate projections to
the whole $Q_{2,3}$ extension.
This proves the common-factor compatibility which
separate $H$ and $T$ comparisons alone would not supply.

For every $\ell\ge5$, the chosen cusp's abelian
classes equal those of $\sigma$ by construction,
so there is a $T_\ell$-conjugator between its
abelian section and that of $\sigma$.
Combine it with the $H_\ell$-conjugator from
(A.102.12).  Equation (A.102.11) allows these
independent pairs and the already constructed
$Q_{2,3}$ conjugator to be assembled into one
geometric element of $Q$.
This proves both cuspidal rows of (A.102.2)
in the actual quotient $\mathcal Q$.

\emph{The other bad places and the explicit good-place failure.}
At two, the altered rational primes $127$ and $23$
are units and their complete local unit groups
have only two-primary components.  The latter
are unchanged, so $A_2=254$, $B_2=-253$ in
the full local completion.
At eleven, the unit group has a pro-eleven part
and root-of-unity group of order ten.
Neither $23$ nor $127$ occurs among these
coefficient primes, so the same conclusion holds.
The $H$ section is already the point section,
and the relative abelian cocycle is supported
only at the independent factors $23,127$ of
(A.102.11).  The local abelian coboundaries can
therefore be removed by geometric elements
$(1,t_\ell)$ in those factors without changing
the $H$ component or $Q_{2,3}$.
This proves the two ordinary rows of (A.102.2).

The displayed Legendre discriminant factors as
\[
       16\,254^2(1-254)^2=2^6\,11^2\,23^2\,127^2.
\tag{A.102.13}
\]
The odd factors have multiplicative reduction
since $c_4$ is a unit there; at two the
contact-one minimal model of A.94 has
discriminant valuation six.
Every other prime has good reduction.
Thus (A.102.2) checks every bad place.

At $47$, reduction of a unit depends on its
exponent modulo $46=2\cdot23$.
The two- and twenty-three-primary coordinates
in (A.102.7)--(A.102.8) give
\[
    \alpha\equiv31\pmod{46},\qquad
    \beta\equiv27\pmod{46},\qquad
    2\,127^{31}\equiv-11\,23^{27}\equiv33\pmod{47}.
\tag{A.102.14}
\]
The sum is $19$ modulo $47$, rather than one.
In the full abelian completion the two valuation
coordinates are zero, so any matching ordinary
point would have both $x$ and $1-x$ units;
its two reductions cannot be $(33,33)$.
The cusp equations are respectively $B=1$,
$A=1$, and $-A/B=1$.  Their reductions here
are $33$, $33$, and $-1$, so none is possible.

The same exclusion already holds at exponent
$46$, without retaining the full valuation.
For a proposed nonintegral ordinary point,
$v_{47}(x)>0$ would give $1-x\equiv1$;
$v_{47}(1-x)>0$ would give $x\equiv1$;
and $v_{47}(x)<0$ would give
$-x/(1-x)\equiv1$ after cancellation of its
equal valuations.
Each contradicts the corresponding residue
$33$, $33$, or $-1$ in (A.102.14).
For the remaining unit-unit case the sum
contradiction applies.  Residues of these
unit classes are well-defined modulo
forty-sixth powers because every such power
has residue one in $\F_{47}^\times$.
Thus the finite geometric abelian quotient
already excludes every point and cusp class.
Any Selmer section lifting $\sigma$ would
have one of those local classes at $47$,
which is impossible.

\emph{Recognition of every labeled integral realization at $254$.}
Now let $s$ be a noncuspidal Selmer section in the
stated realization range.  This paragraph concerns
an actual full section, not the constructed quotient
section $\sigma$.
A.94 gives its exact dyadic valuation pair $(1,0)$
and makes each fixed odd-coefficient global Weil
multiplier a square.
A.16 gives the full global mod-two Kummer classes
\[
                  [A_s]=[254],\qquad [B_s]=[-253]
                              \quad\text{in }\Q^\times/\Q^{\times2}.
\tag{A.102.15}
\]
This includes the signs: $A_s$ has positive sign
and $B_s$ negative sign.  At $127$ the valuation
parities are $(1,0)$, so a locally cuspidal section
must have label zero.  The cusp-one vector has
first valuation zero, and the cusp-infinity vector
has equal valuations; neither has these parities.
At $11$ and $23$ the same argument with parity
$(0,1)$ forces label one.
A.37 excludes every other finite cuspidal place.
At every ordinary odd bad place A.78 gives the
exact local point $254$; at a good odd place both
coordinates of the ordinary point are units.
Together with the exact dyadic pair, these facts
and rational prime factorization in the completion
give
\[
               A_s=2\,127^u,\qquad
               B_s=-11^v23^w,
                           \qquad u,v,w\in\widehat\Z^\times.
\tag{A.102.16}
\]
For the asserted unit property, an ordinary bad
place has the corresponding exponent equal to
one.  At a cusp, the saturated integral Tate
comparison has unit diagonal entries and compares
its exponent with the realized valuation one.
Their ratio is a unit at every coefficient, as
in (A.78.5).
At an odd coefficient $\ell$, its form is
$a_\ell/d_\ell=a_\ell^2/\Delta_\ell$.
A.94 makes $\Delta_\ell$ a square, so all three
odd components $u_\ell,v_\ell,w_\ell$ are unit
squares, whether their physical place is ordinary
or cuspidal.  No square multiplier at two is used;
the two-components are odd because the exponents
are units (also by (A.102.15)).

The place $11$ cannot be cuspidal.  Its only
possible cusp label is one, whose regular
coordinate is $A_s=1$ locally.  At coefficient
eleven this would give
\[
        u_{11}=-\frac{\log_{11}(2)}{\log_{11}(127)}
                       \equiv-1\equiv10\pmod{11}.
\tag{A.102.17}
\]
Both normalized logarithms in this formula
are $6$ modulo eleven by (A.102.5).
The residue $10$ is a nonsquare, contradicting
the unit-square property of $u_{11}$.
Thus $11$ is ordinary and A.78 gives $v=1$
as an equality in the full $\widehat\Z$.

Suppose $s$ were not a point section.
By A.37 its local cusp labels could not be
confined to one label.  After the exclusion of
$11$, the only available labels and primes
are cusp zero at $127$ and cusp one at $23$.
Both must therefore occur.
At $23$ the regular-coordinate equality
$A_s=1$ gives $u_{23}=a_{23}$ exactly.
The full Tate extension at that cusp, with
$B_s=-11\,23^w$, has parameter $B_s^2/256$.
Its valuation is $2w$, whereas the realized
nome has valuation two.  The integral triangular
comparison therefore gives the precise principal-unit
equation
\[
      w_{23}\log_{23}(q_{23}/23^2)
                 =2\log_{23}(11)-8\log_{23}(2).
\tag{A.102.18}
\]
In particular $w_{23}=b_{23}$.
This equation retains the actual extension
class and does not replace an unspecified
linear intertwiner by a symplectic one.
Equations (A.102.7), (A.102.15), and the
odd two-components now force
\[
                 u\equiv31\pmod{46},\qquad
                 w\equiv27\pmod{46}.
\tag{A.102.19}
\]
This conclusion holds regardless of every
other coefficient component of $u,w$.
Transport it through the actual reduction
homomorphism on the rational prime units at
$47$.  It gives precisely the two residues
$(33,33)$ from (A.102.14).
Since $47$ is good, A.37 requires an ordinary
local point, which has already been excluded
by their sum.
This contradiction proves that $s$ is a
point section.  Its ordinary local point
at $11$ is $254$ by A.78, so its rational
parameter is exactly $254$.
\end{proof}

The exact arithmetic certificate is
\path{code/legendre_bad_place_fibre_certificate.py},
with retained output in
\path{code/legendre_bad_place_fibre_certificate_output.txt}.
It computes the inverse nomes and logarithm ratios
modulo $p^7$, lifts all four unit square roots,
checks every bad-place torsion order and the
residues at $47$, and retains the finite search
that selected this witness.
The global quotient section and its simultaneous
local conjugators are constructed in the proof;
no finite calculation is used as a substitute for them.
